\documentclass[10pt,a4paper]{amsart}
\usepackage{amsmath,amssymb,amsthm,graphicx}
% Retained minimal preamble for the Stacks Project chapter varieties.tex.
% Verbatim excerpt of the chapter preamble: Xy-pic setup (native preamble lines 14-18), the native theorem
% declarations with the mathematical macros (native preamble lines 154-217), and the project own no-op
% definitions of the reference/slogan/history display environments (native wrapper.tex lines 1-3) that the
% counted statement uses for its slogan. Copied mechanically; nothing added, removed or redefined.
% For commutative diagrams we use Xy-pic
\usepackage[all]{xy}

% We use 2cell for 2-commutative diagrams.
\xyoption{2cell}
% Theorem environments.
%
\theoremstyle{plain}
\newtheorem{theorem}[subsection]{Theorem}
\newtheorem{proposition}[subsection]{Proposition}
\newtheorem{lemma}[subsection]{Lemma}

\theoremstyle{definition}
\newtheorem{definition}[subsection]{Definition}
\newtheorem{example}[subsection]{Example}
\newtheorem{exercise}[subsection]{Exercise}
\newtheorem{situation}[subsection]{Situation}

\theoremstyle{remark}
\newtheorem{remark}[subsection]{Remark}
\newtheorem{remarks}[subsection]{Remarks}

\numberwithin{equation}{subsection}

% Macros
%
\def\lim{\mathop{\mathrm{lim}}\nolimits}
\def\colim{\mathop{\mathrm{colim}}\nolimits}
\def\Spec{\mathop{\mathrm{Spec}}}
\def\Hom{\mathop{\mathrm{Hom}}\nolimits}
\def\Ext{\mathop{\mathrm{Ext}}\nolimits}
\def\SheafHom{\mathop{\mathcal{H}\!\mathit{om}}\nolimits}
\def\SheafExt{\mathop{\mathcal{E}\!\mathit{xt}}\nolimits}
\def\Sch{\mathit{Sch}}
\def\Mor{\mathop{\mathrm{Mor}}\nolimits}
\def\Ob{\mathop{\mathrm{Ob}}\nolimits}
\def\Sh{\mathop{\mathit{Sh}}\nolimits}
\def\NL{\mathop{N\!L}\nolimits}
\def\CH{\mathop{\mathrm{CH}}\nolimits}
\def\proetale{{pro\text{-}\acute{e}tale}}
\def\etale{{\acute{e}tale}}
\def\QCoh{\mathit{QCoh}}
\def\Ker{\mathop{\mathrm{Ker}}}
\def\Im{\mathop{\mathrm{Im}}}
\def\Coker{\mathop{\mathrm{Coker}}}
\def\Coim{\mathop{\mathrm{Coim}}}

% Boxtimes
%
\DeclareMathSymbol{\boxtimes}{\mathbin}{AMSa}{"02}

%
% Macros for moduli stacks/spaces
%
\def\QCohstack{\mathcal{QC}\!\mathit{oh}}
\def\Cohstack{\mathcal{C}\!\mathit{oh}}
\def\Spacesstack{\mathcal{S}\!\mathit{paces}}
\def\Quotfunctor{\mathrm{Quot}}
\def\Hilbfunctor{\mathrm{Hilb}}
\def\Curvesstack{\mathcal{C}\!\mathit{urves}}
\def\Polarizedstack{\mathcal{P}\!\mathit{olarized}}
\def\Complexesstack{\mathcal{C}\!\mathit{omplexes}}
% \Pic is the operator that assigns to X its picard group, usage \Pic(X)
% \Picardstack_{X/B} denotes the Picard stack of X over B
% \Picardfunctor_{X/B} denotes the Picard functor of X over B
\def\Pic{\mathop{\mathrm{Pic}}\nolimits}
\def\Picardstack{\mathcal{P}\!\mathit{ic}}
\def\Picardfunctor{\mathrm{Pic}}
\def\Deformationcategory{\mathcal{D}\!\mathit{ef}}
\newenvironment{reference}{}{}
\newenvironment{slogan}{}{}
\newenvironment{history}{}{}
\usepackage{cleveref}
\crefname{originallemma}{original lemma}{original lemmas}
\Crefname{originallemma}{Original lemma}{Original lemmas}
\makeatletter
\newlabel{algebra-lemma-spec-product}{{00ED}{source}{}{sourcepointer-1}{}}
\newlabel{algebra-lemma-spec-product@cref}{{[originallemma][0][]00ED}{[1][0][]source}{}{}{}}
\makeatother
\crefname{originallemma}{original lemma}{original lemmas}
\Crefname{originallemma}{Original lemma}{Original lemmas}
\makeatletter
\newlabel{coherent-lemma-flat-base-change-cohomology}{{02KH}{source}{}{sourcepointer-2}{}}
\newlabel{coherent-lemma-flat-base-change-cohomology@cref}{{[originallemma][0][]02KH}{[1][0][]source}{}{}{}}
\makeatother
\crefname{originallemma}{original lemma}{original lemmas}
\Crefname{originallemma}{Original lemma}{Original lemmas}
\makeatletter
\newlabel{lemma-proper-geometrically-reduced-global-sections}{{0BUG}{source}{}{sourcepointer-3}{}}
\newlabel{lemma-proper-geometrically-reduced-global-sections@cref}{{[originallemma][0][]0BUG}{[1][0][]source}{}{}{}}
\makeatother
\crefname{originallemma}{original lemma}{original lemmas}
\Crefname{originallemma}{Original lemma}{Original lemmas}
\makeatletter
\newlabel{schemes-lemma-morphism-into-affine}{{01I1}{source}{}{sourcepointer-4}{}}
\newlabel{schemes-lemma-morphism-into-affine@cref}{{[originallemma][0][]01I1}{[1][0][]source}{}{}{}}
\makeatother
\title{The fibres of the canonical morphism to the spectrum of global functions are geometrically connected}
\author{The Stacks Project Authors}
\date{stacks-project:varieties.tex:local-snapshot-20261011; GNU FDL 1.2 or later}
\begin{document}
\maketitle
\noindent\textit{Editorial note.} All delivered mathematical text is the author's own, copied from the pinned Stacks Project chapter file varieties.tex; no statement or proof step was rewritten, reordered or removed.\par\medskip
\noindent\textit{Editorial note.} The proof is complete as written and ends with the contradiction in the local ring A\_i; no step is referred to the reader.\par\medskip
\noindent\textit{Editorial note.} The counted spans carry 5 navigation stubs and no bibliography citation.\par\medskip
\noindent\textit{Editorial note.} The preamble is the chapter's retained minimal preamble (sources/stacks-varieties/retained-preamble.tex lines 1-77), a verbatim excerpt of the native preamble.\par\medskip
\noindent\textit{Editorial note.} The label lemma-baby-stein is defined only in varieties.tex, so no cross-chapter restatement had to be weighed for this card.\par\medskip
\noindent\textit{Allowed standard prerequisites (editorial).}\begin{itemize}
\item Global sections of a proper geometrically reduced scheme over a field form a finite product of local Artinian algebras. Same chapter, Lemma lemma-proper-geometrically-reduced-global-sections, invoked at proof lines 2024-2026 and again at 2062-2064. Supplies the decomposition A = product of A\_i and the identification A\_i = Gamma(X\_i, O\_\{X\_i\}).
\item The canonical morphism to the spectrum of the global sections. Schemes, Lemma schemes-lemma-morphism-into-affine, invoked at proof lines 2020-2023. Supplies the identification of the morphism X to Spec(A) as the one induced by the ring A.
\item Flat base change for the cohomology of the structure sheaf. Cohomology of Schemes, Lemma coherent-lemma-flat-base-change-cohomology, invoked at proof lines 2043-2046. Supplies Gamma(X-bar, O) = A-bar after base change to the algebraic closure, which allows the reduction to k algebraically closed.
\item The product decomposition of the spectrum corresponds to the product decomposition of the ring. Algebra, Lemma algebra-lemma-spec-product, invoked at proof lines 2057-2060. Supplies Spec(A) = disjoint union of the Spec(A\_i).
\end{itemize}\medskip
\ifdefined\hypertarget\hypertarget{sourcepointer-1}{}\fi
\noindent\textit{Source reference.} Original lemma 00ED: sources/stacks-chow/tags.txt, source lines 530--530. Pointer outside the counted span, invoked once for context; the load-bearing argument of the counted statement is the author's proof reproduced in full in the proof span.\par\medskip
\ifdefined\hypertarget\hypertarget{sourcepointer-2}{}\fi
\noindent\textit{Source reference.} Original lemma 02KH: sources/stacks-chow/tags.txt, source lines 3372--3372. Pointer outside the counted span, invoked once for context; the load-bearing argument of the counted statement is the author's proof reproduced in full in the proof span.\par\medskip
\ifdefined\hypertarget\hypertarget{sourcepointer-3}{}\fi
\noindent\textit{Source reference.} Original lemma 0BUG: sources/stacks-varieties/pinned-original-source.tex, source lines 1933--1933. Pointer outside the counted span, invoked once for context; the load-bearing argument of the counted statement is the author's proof reproduced in full in the proof span.\par\medskip
\ifdefined\hypertarget\hypertarget{sourcepointer-4}{}\fi
\noindent\textit{Source reference.} Original lemma 01I1: sources/stacks-chow/tags.txt, source lines 1962--1962. Pointer outside the counted span, invoked once for context; the load-bearing argument of the counted statement is the author's proof reproduced in full in the proof span.\par\medskip
\begin{lemma}
\label{lemma-baby-stein}
Let $X$ be a proper scheme over a field $k$. Set
$A = H^0(X, \mathcal{O}_X)$. The fibres of the canonical
morphism $X \to \Spec(A)$ are geometrically connected.
\end{lemma}

\begin{proof}
Set $S = \Spec(A)$. The canonical morphism $X \to S$
is the morphism corresponding to
$\Gamma(S, \mathcal{O}_S) = A = \Gamma(X, \mathcal{O}_X)$ via
Schemes, Lemma \ref{schemes-lemma-morphism-into-affine}.
The $k$-algebra $A$ is a finite product $A = \prod A_i$
of local Artinian $k$-algebras finite over $k$, see
Lemma \ref{lemma-proper-geometrically-reduced-global-sections}.
Denote $s_i \in S$ the point corresponding to the maximal
ideal of $A_i$. Choose an algebraic closure $\overline{k}$ of $k$ and
set $\overline{A} = A \otimes_k \overline{k}$.
Choose an embedding $\kappa(s_i) \to \overline{k}$ over $k$; this
determines a $\overline{k}$-algebra map
$$
\sigma_i : \overline{A} = A \otimes_k \overline{k} \to
\kappa(s_i) \otimes_k \overline{k} \to \overline{k}
$$
Consider the base change
$$
\xymatrix{
\overline{X} \ar[r] \ar[d] & X \ar[d] \\
\overline{S} \ar[r] & S
}
$$
of $X$ to $\overline{S} = \Spec(\overline{A})$. By
Cohomology of Schemes, Lemma
\ref{coherent-lemma-flat-base-change-cohomology} we have
$\Gamma(\overline{X}, \mathcal{O}_{\overline{X}}) = \overline{A}$.
If $\overline{s}_i \in \Spec(\overline{A})$ denotes the
$\overline{k}$-rational point corresponding to $\sigma_i$,
then we see that $\overline{s}_i$ maps to $s_i \in S$
and $\overline{X}_{\overline{s}_i}$ is the base change of
$X_{s_i}$ by $\Spec(\sigma_i)$. Thus we see that it suffices to prove
the lemma in case $k$ is algebraically closed.

\medskip\noindent
Assume $k$ is algebraically closed. In this case $\kappa(s_i)$
is algebraically closed and we have to show that $X_{s_i}$
is connected. The product decomposition
$A = \prod A_i$ corresponds to a disjoint union decomposition
$\Spec(A) = \coprod \Spec(A_i)$, see
Algebra, Lemma \ref{algebra-lemma-spec-product}.
Denote $X_i$ the inverse image of $\Spec(A_i)$.
It follows from
Lemma \ref{lemma-proper-geometrically-reduced-global-sections} part (2) that
$A_i = \Gamma(X_i, \mathcal{O}_{X_i})$.
Observe that $X_{s_i} \to X_i$ is a closed immersion inducing
an isomorphism on underlying topological spaces (because $\Spec(A_i)$
is a singleton). Hence if $X_{s_i}$ isn't connected, then neither is
$X_i$. So either $X_i$ is empty and $A_i = 0$ or $X_i$ can be written
as $U \amalg V$ with $U$ and $V$ open and nonempty which would
imply that $A_i$ has a nontrivial idempotent. Since $A_i$ is
local this is a contradiction and the proof is complete.
\end{proof}

\end{document}
